\documentclass{minmal}
\begin{document}
\textbf{Arbeidsopgaver fra Loops:}\\\\
1. Forklar med ord hva som menes med begrepet spesifikk varmekapasitet. Hva er benevningen? \\\\ Spesifikk vafmekapsitet sier hvor mye energi som trengs for å varme opp 1kg av stoffet $1^\circ C$.  \\\\
2. Hvilken verdi har vannets spesifikke varmekapasitet ($C_p$)?  \\\\ Vannets spesfikke varmekapasistet er $C_{p,vann}=4,2kJ/kg^\circ C$.\\\\
3. Du skal varme 3,2 kg. vann fra 34,50C til 67,80C. Hvor mye energi trengs til dette? (svar: 447,552kJ)\\\\
Her gjelder følgende formel:\\\\
$Q=m\cdot C_p\cdot \Delta T$\\\\\
Der:\\
Q=energimengfden (kJ)\\
m=massen (kg)\\
$C_p$=spesifikk varmekapasitet ($kJ/kg^0 C)$\\
$\Delta T$=endring i temperatur ($^\circ C)$\\\\
Vi setter inn verdier:\\
$Q=m\cdot C_p\cdot \Delta T=3,2kg\cdot 4,2kJ/kg^0C\cdot (67,8-34,5)^0C=447,6kJ$
\\\\\
4. Forklar hva menes med begrepet spesifikk smeltevarme. Hva er benevningen? \\\\
Spesifikk smletvarme er en energien som trengs for å smelte 1kg av et stoff.  Enehet: kJ/kg.
5.  Hvilken verdi har vannets spesifikke smeltevarme ($l_s$) ? \\\\
Svar:  $l_{s,vann}=334kJ/kg$
6.  Hvor mye energi trengs for å smelte 4,56 kg is? (svar:1523,04kJ) \\\\
Svar: $Q_s=m\cdot l_s=4,56kg\cdot 334kJ/kg=1523,04kJ$
\newpage\noindent
7.  Forklar hva som menes med begrepet spesifikk fordampingsvarme. Hva er benevningen? \\\\
Spesifkk fordamoingsvarme er den energien som trengs for å fordampe 1kg av stoffet.  Enhet: kJ/kg.\\\\
8. Hvilken verdi har vannets spesifikke fordampingsvarme ($l_f$) ?\\\\  $l_{f,vann}=2256kJ/kg$.\\\\
9. Hvor mye energi trengs for å fordampe 4,56 kg. vann? (svar: 10287,36kJ) \\\\
Formel: $Q_{f,vann}=m\cdot l_{f,vann}=4,56kg\cdot 2256kJ/kg=10287,36kJ$\\\\
10. Hvor stor energi trengs for å fordampe 3 kg is med temperatur $-12^\circ C$? (svar: 9105,6kJ) \\\\
Denne prosessen skjer i 4 trinn:\\\\
1. Vi må varme opp isen til $0^0C$.  Is har en spesifikk varmekapasitet lik $C_{p,is}=2,1kJ/kg^0C$.  Energimengden for på varme opp 3kg is fra $-12^0C$ til $0^0C$:\\\\
$Q_1=m\cdot C_{p,is}\cdot \Delta T=3kg\cdot 2,1kJ/kg^0C\cdot (0--12)^0C=75,6kJK$\\\\
2. Så må vi smelte isen:\\\\
$Q_2=m\cdot l_{s,is}=3kg\cdot 334kJ/kg=1002kJ$\\\\
3. Så må vi varme opp vannet til $100^0C$:\\\\
$Q_3=m\cdot C_{p,vann}\Delta T=3kg\cdot 4,2kJ/kg^0C\cdot 100^0C=1260kJ$\\\\
4. Tilslutt må vi fordampe vannet:\\\\
$Q_4=m\cdot l_{f,vann}=3kg\cdot 2256kJ/kg=6768kJ$\\\\
Energi totalt:
$Q_{totalt}=Q_1+Q_2+Q_3+Q_4=75,6kJ+1002kJ+1260kJ+6768kJ=9105,6kJ$
\newpage

11. Hvor stor del av den totale energien er fordampingen? (svar: 74,3\%)\\\\Energien som brukes til å fordampe 3 kg vann er $Q_4$ og den er lik  6768kJ.   Den totale energein er 9105,6kJ.  Prosentandelen blir da:\\\\

$$\frac{6768kJ\cdot 100\%}{9105,6kJ}=74,3\%$$

12. Hvor mange kg vann kan du varme $50^\circ C (\Delta T=50^\circ C$) når du har 5 kg. damp som kondenserer? (Svar: ? kg. vann)\\\\
Vi må finne den energimengden som blir avgitt når 5kg damp kondenserer:\\
$$Q_f=m\cdot L_f=2256kJ/kg\cdot 5kg=11280kJ$$\\
Energimegnden som vannet tar opp følger av denne ligningen:\\
$$Q=m\cdot C_p\cdot \Delta T$$\\
Her kjenner vi Q=11280kJ, $C_{P,vann}=4,2kJ/kg^oC\ \Delta T=50^0C$\\
løser m.h.p.  m:
$$m=\frac{Q}{C_{P,vann}\cdot \Delta T}=\frac{11280kJ}{4,2kJ/kg^0C\cdot 50^0C }=53,7kg$$
\newpage
\begin{center}
    
{\textbf{Arbeidsoppgaver:}}
\end{center}
Gitt følende varmeveklser:\\
\varmeveksler{222,2}{3}{?}{78}{2}{5}
Bereng temperatruen ut på kald side.\newpage

Vannet fra Porsnes bassenget skal ha en temperatur på $T_{V_ut}=37^0C.$  Hvor mye kjøevann trenger vi?\\
Vi får avløpsvanmn fra Porsnesbassenget ( 87kg/s) og vannet har en temperatur på $T_{v,inn}=47^0C$.  Dette vannet kjøles med vann fra Tista som ha en temperatur på $T_{K,inn}=10^0c.$  Vanne fra Porsnes bassenget skal ha en temperatur på $T_{V_ut}=37^0C.$  Hvor mye kjøLevann trenger vi?\\
\varmeveksler{47}{37}{23}{10}{87}{?}\newpage
Gitt følgende vaermeveklser:
\varmeveksler{90}{75}{45}{?}{20}{20}
Hva er temperaturen inn på kald side ($T_{k,inn}$)?\nl[2]
Hvor mye energi trengs for å varme opp 5,6kg vann fra $20^\circ C$ til $40^\circ C$?\nl[2]
Hvor mye energi trengs for å smelte 3,7kg is som har en temperatur på $0^\circ C$?\nl[2]
Hvor mye energi trengs for å smelte 3,7kg is som har en temperatur på $-20^\circ C$?\nl[2]
Hvor mye energi trengs for å fordampe 3,7kg is som har en temperatur på $-20^\circ C$?\newpage
\noindent\textbf{Viktige data og formler:}\nl[2]

\noindent Vannets spesifikke varmekapasitet  ($C_{P,vann}=4,2kJ/kg^\circ C$\\
Is sin spesifikke varmekapasitet  ($C_{P,is}=2,1kJ/kg^\circ C$\\
Vannets spesifikke fordampingsvarme  ($l_f=2256kJ/kg$)\\
Vannets spesifikke smeltevarme  ($l_s=334kJ/kg$)\nl[2]

\noindent\textbf{Beregning av fordampingsvarme:}\\
$Q_f=m\cdot l_f$\\\\
$Q_f$:fordampingsvarme ($kJ)$
$m$:masse av stofffet ($kg$)
$l_f$:Spesifikk fordampingsvarme ($2256kJ/kg$)\nl[2]

\noindent\textbf{Beregning av smeltevarme:}\\
$Q_s=m\cdot l_s$\\\\
$Q_s$:smeltevarme ($kJ)$\\
$m$:masse av stoffet ($kg$)\\
$l_s$:Spesifikk smeltevarme ($334kJ/kg$)\nl[2]
\noindent\textbf{Beregning av energi for å varme stoffet:}\\
$Q=m\cdot C_{P}\cdot \Delta T$\\\\
$m$:masse av stoffet ($kg$)\\
$C_P$: Spesifik varmekapasitet ($kJ/kg^\circ C$)
$\Delta T$: temperaturforskjell ($^\circ C$)
\newpage
\varmeveksler{73}{53}{30}{20}{30}{?}
Vi begynner med å beregne den energien som den varme væsken avgir:\\
$$Q_v=F_v\cdot C_P\cdot \Delta T=30kg/s\cdot 4,2kJ/kgC\cdot (73-53)^\circ C=2520kW$$
Viktig poeng nå er at dette er den energien den kalde væsken vil ta opp.   Om vi løser lingnen m.h.p. F fåer vi om vi bruker trekanten:\\
$$F=\frac{Q_k}{C_P\cdot \Delta T}=\frac{2520kW}{4,2kJ/kgC\cdot (20-10)^\circ C}=60kg/s$$
\newpage
Gitt følgende varmeveksler:\\
\varmeveksler{80}{?}{15}{30}{30}{20}
Her kjenner vi alt på kald side:\\
$$Q_k=F\cdot C_P\cdot \Delta T=20kg/s\cdot 4,2kJ/kg^\circ C\cdot (30-15)^\circ C=1260kW $$
Nå bruker vi trekanten for må finne  $\Delta T$:
$$\Delta T=\frac{Q}{F\cdot C_P}=\frac{1260kW}{30kg/s\cdot 4,2kJ/kg^\circ C}=10^\circ C$$
Det vi har funnet ut nå er at den varme væsken vil endre temperaturen med $10^\circ C$.  Den varme væsken vil selvsagt bli avkjølt.  Vi får:
$$T_{v,ut}=T_{v, inn}-\Delta T=80^\circ C-10^\circ C=70^\circ C$$
\newpage
Beregn $F_k$:
\varmeveksler{60}{30}{25}{10}{20}{?}
Vi beregner først energien på varm side ($Q_v$):
$$Q_v=F_v\cdot C_p\cdot \Delta T=20kg/s\cdot 4,2kJ/kg^\circ \cdot(60-30)^\circ C=2520kW$$
For $F$ kan vi skrive:
$$F=\frac{Q_c}{C_P\cdot \Delta T=}=\frac{2520kW}{4,2kJ/lkg^\circ C\cdot(25-10)^\circ C}=40kg/s$$


\newpage
\noindent \textbf{Viktige spørsmål til neste prøve:}\\\\
1.  definer begrepene spesifikk varmekapasitet, spesifikk fordampingsvarme og spesifikk smeltevarme.\nl[2]
2.  Finn verdien for disse for vann.\nl[2]
3. Hvor mye enerrgi trengs for å varme opp 3,5kg vann fra $20^\circ  C$ til $45^\circ C$?\nl[2]
4. Hvor mye energi trengs for å smelte 5,6kg is som har  en temperatur på $0^\circ  C$?\nl[2]
5. Hvor mye energi trengs for å fordampe 5,6kg vann som har  en temperatur på $100^\circ  C$?\nl[2]\newpage


6. Hvor mye energi trengs for å fordampe 7,8kg is som har en temperatur på $-7,6^\circ C$?  Hint: Is har en spesifikk varmekapasitet på $2,1kJ/kg^\circ C$\\
Her har vi 4 ulike prosesser.  \\
1. Vi må varme opp isen fra $-7,6^\circ C$ til $0^\circ C$  (da er  $C_p=2,1kJ/kg^\circ  C$\\
2. Vi må smelte isen (is overføres til vann)\\
3. varme til kokepunktet.\\
4. fordampe vannet. \\
Vi beregner hvert trinn:\\
trinn 1:  
$$Q_1=m\cdot C_{p,is}\cdot \Delta T=7,8kg\cdot 2,1kJ7kg^\circ  C\cdot 7,6^\circ C=124,5kJ $$\\
Trinn 2: Vi smelter isen:
$$Q_S=m\cdot L_s=7,8kg\cdot 334kJ/kg=2605,2kJ$$\\
Trinn 3: Vi skal varme opp vannet til kokepunkt $100^\circ C$: \\
$$Q_2=m\cdot C_{p,vann}\cdot \Delta T=7,8kg\cdot 4,2kJ/kg^\circ C\cdot 100^\circ C=3276kJ$$\\
Trinn 4: Vi fordamper vannet:\\
$$Q_f=m\cdot L_f=7,8kg\cdot 2256kJ/kg=17596,8kJ$$
Total enegimengde: $124,5kJ+2605,2kJ+3276kJ+17596,8kJ=23602,5kJ$\\
Når noe nav dere har funent ut hvor stor prosentandelen fordampingen er av det totale,  så er det game over.

8. Beregn temperaturen inn på kald side i denne varmeveksleren (svar: $T_{k,inn}=14,33^\circ C$).\\
\varmeveksler{83}{67}{25}{?}{10}{15}
9.  Hor stor er flowen på varm side? (svar $F_v=17,14kg/s$).\nl[2]
\varmeveksler{67}{53}{12 }{32 }{?}{12}
\newpage
\varmeveksler{100}{100}{? }{10 }{1}{12}
\end{document}